\chapter{Additional Figures, Evidence Audit, and Remaining Requirements}


\section{Cross-condition result map}
The nodes summarize verified NDCG@10 directions for three-way minus two-way. This is a categorical evidence map; it has no numerical axes and does not encode confiden\allowbreak{}ce-inter\allowbreak{}val widths. The BCA node refers to its primary small-component view, rather than its combined descriptive average.

\begin{figure}[p]
\centering
\begin{tikzpicture}[>=Latex, every node/.style={draw=black!65, rounded corners=2pt, fill=black!3, text width=4.35cm, minimum height=0.85cm, align=center, font=\small}, every path/.style={draw=black!60, thick}]
\node (A) at (0.00,0.00) {Add pretrained E5 to BM25 plus BGE-M3};
\node (B) at (-5.05,-1.55) {ALQAC clean: small negative difference};
\node (C) at (0.00,-1.55) {ALQAC parsed: small negative difference};
\node (D) at (5.05,-1.55) {Zalo clean: small positive difference};
\node (E) at (-2.52,-3.10) {Zalo parsed: positive paired evidence};
\node (F) at (2.52,-3.10) {BCA small components: no established benefit};
\draw[->] (A.south) -- (B.north);
\draw[->] (A.south) -- (C.north);
\draw[->] (A.south) -- (D.north);
\draw[->] (A.east) -- (7.5,0.00) |- (E.east);
\draw[->] (A.west) -- (-7.5,0.00) |- (F.west);
\end{tikzpicture}
\caption[Direction of the incremental E5 result by evaluation condition.]{Direction of the incremental E5 result by evaluation condition.}
\label{fig:G.1}
\end{figure}
The nodes compare the direction of fine-tuned minus pretrained NDCG@10 against each cohort's own baseline. No arrow represents a paired numerical subtraction between development and held-out cohorts. The diagram's purpose is to expose the reversal that would be hidden by reporting the selected pilot alone.

\begin{figure}[p]
\centering
\begin{tikzpicture}[>=Latex, every node/.style={draw=black!65, rounded corners=2pt, fill=black!3, text width=4.35cm, minimum height=0.85cm, align=center, font=\small}, every path/.style={draw=black!60, thick}]
\node (A) at (0.00,0.00) {Development selection: both hybrids improve};
\node (B) at (0.00,-1.55) {Lock recipe and train three seeds};
\node (C) at (0.00,-3.10) {Corrected held-out comparison};
\node (D) at (-5.05,-4.65) {Two-way: all three NDCG differences negative};
\node (E) at (0.00,-4.65) {Three-way: all three NDCG differences negative};
\node (F) at (5.05,-4.65) {Cluster intervals span zero};
\draw[->] (A.south) -- (B.north);
\draw[->] (B.south) -- (C.north);
\draw[->] (C.south) -- (D.north);
\draw[->] (C.south) -- (E.north);
\draw[->] (C.south) -- (F.north);
\end{tikzpicture}
\caption[Development and corrected held-out outcome of hybrid adaptation.]{Development and corrected held-out outcome of hybrid adaptation.}
\label{fig:G.2}
\end{figure}
